\documentclass[11pt,a4paper]{article}

\usepackage[T1]{fontenc}
\usepackage[utf8]{inputenc}
\usepackage{lmodern}
\usepackage[margin=2cm]{geometry}
\usepackage{xcolor}
\usepackage[colorlinks=true,urlcolor=blue!60!black]{hyperref}

\setlength{\parindent}{0pt}
\setlength{\parskip}{5pt}
\pagestyle{empty}

\begin{document}

\hfill\begin{tabular}{l}
Instituto de Astronom\'ia Te\'orica y Experimental (IATE)\\
CONICET -- Universidad Nacional de C\'ordoba\\
Laprida 854, X5000BGR C\'ordoba, Argentina\\
\href{mailto:fdavilakurban@unc.edu.ar}{fdavilakurban@unc.edu.ar}\\[8pt]
November 2026
\end{tabular}

\vspace{4pt}
Prof.~Uro\v{s} Seljak, Director, and the BCCP Postdoctoral Fellowship Selection Committee\\
Berkeley Center for Cosmological Physics\\
Department of Physics, University of California, Berkeley\\
Berkeley, CA 94720, USA

\vspace{6pt}
Dear Prof.~Seljak and members of the Selection Committee,

I am applying for a BCCP Postdoctoral Fellowship starting in 2027. I received my PhD in Astronomy from the Universidad Nacional de C\'ordoba in 2022, and I currently carry out research in large-scale structure statistics at the Instituto de Astronom\'ia Te\'orica y Experimental (IATE) in C\'ordoba, Argentina, where I also hold an Assistant Professor position at the university. At BCCP I propose to test whether the geometry of the cosmic web, and not only its density, shapes the baryon acoustic feature, and whether DESI can use that information to sharpen or to safeguard its distance measurements.

The proposal grows out of the project I currently lead at IATE on the dependence of the BAO signal on cosmic-web environment. I designed it and built its pipeline, from DisPerSE filaments and density-matched random dilutions to analytic covariances and \texttt{desilike} BAO fits. The first paper, on MultiDark, the Flagship light cone and SDSS, is in preparation, and preliminary measurements give a first hint of a narrower acoustic peak near filaments. Density-based studies of environment and BAO average over directions, yet the flows that smear the peak are anisotropic around filaments. At BCCP I would measure the acoustic feature as a function of the angle between pair separations and the local filament axis, test whether standard reconstruction absorbs that dependence, and carry the measurement to DESI. Because reconstruction concentrates the methodological systematics of BAO, either outcome tests the step of the analysis that matters most.

BCCP is the natural place for this work. DESI is led from the Lawrence Berkeley National Laboratory, and the project connects to work by Prof.~Martin White, Dr.~Joanne Cohn and Dr.~Pat McDonald on the theory of reconstruction and on the BAO modelling systematics of DESI. A second line that I lead, on the void probability function as a probe of primordial non-Gaussianity, connects to the DESI measurement of $f_{\rm NL}$ by Dr.~Edmond Chaussidon and collaborators. 

In return, I bring tools that would serve the galaxy-clustering projects of BCCP members, including in the grant-supported positions for which this call also considers applicants: BAO fitting and covariance estimation, cosmic-web and void identification, higher-order clustering statistics, and efficient correlation-function estimators, including a glass-catalogue estimator developed with Dr.~Ariel S\'anchez that cuts the dominant pair-counting cost by a factor of about 86. I have led most of my projects as first author, most recently a cross-correlation of nearby galaxies with Pierre Auger cosmic-ray data (under review), and I have co-organised three meetings at IATE, including the 2024 conference \textit{Mind the gap: Galaxies and the Large Scale Structure}.

I attach my CV, publication list and statement of research interests; three referees will submit their letters through AcademicJobsOnline. Thank you for considering my application.

\vspace{4pt}
\hspace{0.45\textwidth}Sincerely,

\vspace{16pt}
\hspace{0.45\textwidth}Dr.\ Federico D\'avila-Kurb\'an

\end{document}
